\chapter*{Annexes Techniques}
\label{chap:appendix}
\markboth{\MakeUppercase{Annexes Techniques}}{}
\addcontentsline{toc}{chapter}{Annexes Techniques}

\section*{Annexe 1 : Manifeste du Module (\texttt{\_\_manifest\_\_.py})}
Le fichier de déclaration du module spécifie les métadonnées pour \textbf{\companyName}, les dépendances vers les modules natifs d'Odoo 18, ainsi que l'ordonnancement précis de chargement des vues, modèles multi-sites, assistants modaux et règles de sécurité.

\begin{lstlisting}[language=Python, caption={Extrait du fichier de déclaration \texttt{\_\_manifest\_\_.py}}]
{
    'name': 'Equipment Maintenance System',
    'version': '18.0.1.0.0',
    'category': 'Manufacturing/Maintenance',
    'summary': 'Industrial Maintenance, Multi-Site Locations, Analytics & BYO-LLM Assistant',
    'sequence': 10,
    'license': 'LGPL-3',
    'author': 'IT-Koncept SA & Engineering Team',
    'website': 'https://www.it-koncept.ch',
    'depends': [
        'base',
        'mail',
        'stock',
        'product',
        'web',
    ],
    'data': [
        # Securite et regles d'acces
        'security/equipment_security.xml',
        'security/ir.model.access.csv',
        # Sequences et Crons
        'data/equipment_sequence.xml',
        'data/equipment_cron.xml',
        'data/mail_template_data.xml',
        'data/checklist_template_data.xml',
        # Rapports QWeb PDF
        'report/equipment_report_templates.xml',
        'report/equipment_reports.xml',
        # Vues Metier & Navigation
        'views/equipment_location_views.xml',
        'views/equipment_category_views.xml',
        'views/equipment_views.xml',
        'views/maintenance_order_views.xml',
        'views/repair_request_views.xml',
        'views/equipment_template_views.xml',
        'views/equipment_dashboard_views.xml',
        'views/equipment_ai_config_views.xml',
        'views/equipment_menu.xml',
        # Wizards modaux
        'wizard/equipment_ai_nl_wizard_views.xml',
        'wizard/equipment_report_wizard_views.xml',
    ],
    'demo': [
        'demo/equipment_demo.xml',
    ],
    'installable': True,
    'application': True,
    'auto_install': False,
}
\end{lstlisting}

\newpage
\section*{Annexe 2 : Matrice de Sécurité \& Règle de Lecture Corrigée}
L'extrait ci-dessous illustre la définition des permissions CRUD dans \texttt{ir.model.access.csv} ainsi que la correction de la règle d'enregistrement \textit{write-own} où \texttt{perm\_read="0"} a été explicitement défini pour empêcher le masquage des lignes de pièces détachées pour l'administrateur.

\begin{lstlisting}[caption={Extrait des permissions CSV et de la règle XML corrigée}]
# ir.model.access.csv (extrait sur 45 regles)
id,name,model_id:id,group_id:id,perm_read,perm_write,perm_create,perm_unlink
access_equipment_viewer,access.equipment.viewer,model_equipment_equipment,group_equipment_viewer,1,0,0,0
access_equipment_tech,access.equipment.tech,model_equipment_equipment,group_equipment_tech,1,1,1,0
access_equipment_mgr,access.equipment.mgr,model_equipment_equipment,group_equipment_manager,1,1,1,1
access_location_user,access.location.user,model_equipment_location,group_equipment_viewer,1,0,0,0
access_location_mgr,access.location.mgr,model_equipment_location,group_equipment_manager,1,1,1,1
access_ai_config_mgr,access.ai.config.mgr,model_equipment_ai_config,group_equipment_manager,1,1,1,1
access_ai_pred_mgr,access.ai.pred.mgr,model_equipment_ai_prediction,group_equipment_manager,1,1,1,0

<!-- Correction de la regle write-own dans equipment_security.xml -->
<record id="rule_spare_part_line_technician_write" model="ir.rule">
    <field name="name">Spare Parts: technician can only modify own lines</field>
    <field name="model_id" ref="model_equipment_spare_part"/>
    <field name="groups" eval="[(4, ref('group_equipment_tech'))]"/>
    <field name="perm_read" eval="0"/>  <!-- Correction clef: lecture non filtree -->
    <field name="perm_write" eval="1"/>
    <field name="perm_create" eval="1"/>
    <field name="perm_unlink" eval="1"/>
    <field name="domain_force">[('create_uid', '=', user.id)]</field>
</record>
\end{lstlisting}

\section*{Annexe 3 : Extrait de la Suite de Tests Backend Odoo (138 Tests)}
L'extrait suivant présente la validation unitaire du calcul du score de santé, de la prévention contre les divisions par zéro sur le MTBF, et du cycle de vie des ordres.

\begin{lstlisting}[language=Python, caption={Extrait des tests unitaires backend (\texttt{tests/test\_equipment.py})}]
from odoo.tests.common import TransactionCase
from odoo.exceptions import ValidationError

class TestEquipmentMaintenance(TransactionCase):

    @classmethod
    def setUpClass(cls):
        super().setUpClass()
        cls.category = cls.env['equipment.category'].create({
            'name': 'Usinage Numerique',
            'code': 'CNC-5AX',
        })
        cls.location = cls.env['equipment.location'].create({
            'name': 'Atelier Production A',
            'location_type': 'room',
        })
        cls.equipment = cls.env['equipment.equipment'].create({
            'name': 'Centre d Usinage 5-Axes',
            'serial_no': 'CH-NYON-2026-01',
            'category_id': cls.category.id,
            'location_id': cls.location.id,
            'purchase_date': '2025-01-01',
            'warranty_end_date': '2027-01-01',
        })

    def test_health_score_new_machine(self):
        """Verifie le score initial optimal d'un equipement neuf"""
        self.equipment._compute_health_score()
        self.assertGreaterEqual(self.equipment.health_score, 90.0)

    def test_mtbf_zero_division_guard(self):
        """Verifie que l'absence de panne ne declenche aucune division par zero"""
        risk_score = self.env['equipment.ai.service'].calculate_failure_risk(self.equipment.id)
        self.assertGreaterEqual(risk_score, 0.0)
        self.assertLessEqual(risk_score, 100.0)
\end{lstlisting}

\newpage
\section*{Annexe 4 : Extrait des Tests E2E de Navigation Playwright}
L'extrait ci-dessous, issu de la suite Playwright (\texttt{tests/e2e/full-demo.spec.ts}), démontre l'automatisation de la navigation : dépliage des listes groupées Odoo 18, contrôle du fil d'Ariane, saisie dans le wizard modal et validation de la confirmation d'écriture.

\begin{lstlisting}[language=Python, caption={Extrait du scénario de navigation Playwright E2E (\texttt{full-demo.spec.ts})}]
import { test, expect } from '@playwright/test';

test.describe('Navigation & View Testing Suite', () => {

  test('Navigate through Equipment list, grouped expansion and detail form', async ({ page }) => {
    // Connexion et navigation vers le menu Equipements
    await page.goto('/web#action=equipment_maintenance.action_equipment_equipment');
    await page.waitForSelector('.o_list_table');

    // Gestion du comportement specifique Odoo 18: deplier les groupes
    const groupHeaders = page.locator('tr.o_group_header');
    if (await groupHeaders.count() > 0) {
      await groupHeaders.first().click();
      await page.waitForTimeout(300);
    }

    // Clic sur la premiere fiche et verification du fil d'Ariane
    await page.locator('.o_data_row').first().click();
    await expect(page.locator('.breadcrumb-item.active')).toBeVisible();

    // Verifier le badge d'indice de sante
    const healthBadge = page.locator('field[name="health_score"]');
    await expect(healthBadge).toBeVisible();
  });

  test('Validate NL Wizard modal navigation without premature close', async ({ page }) => {
    // Ouverture de l'assistant de journalisation en langage naturel
    await page.locator('button:has-text("Journaliser une intervention (IA)")').click();
    await expect(page.locator('.modal-dialog')).toBeVisible();

    // Saisie du rapport d'intervention
    await page.fill('textarea[name="raw_text"]', 'Remplacement filtre a huile sur CNC-5AX suite vibration');
    await page.click('button:has-text("Analyser le texte")');

    // Assertion critique: le dialogue doit rester ouvert (reload action) avec apercu
    await expect(page.locator('.modal-dialog')).toBeVisible();
    await expect(page.locator('field[name="extracted_summary"]')).not.toBeEmpty();
  });

});
\end{lstlisting}

\section*{Annexe 5 : Script d'Orchestration E2E sous Docker/WSL (\texttt{scripts/e2e-run.sh})}
Ce script PowerShell/Bash assure la gestion du cycle de vie du conteneur Docker Odoo 18, la persistance des volumes et le déclenchement de la suite Playwright.

\begin{lstlisting}[caption={Extrait du script d'exécution \texttt{scripts/e2e-run.sh}}]
#!/usr/bin/env bash
set -euo pipefail

COMPOSE_FILE="docker-compose.yml"
TARGET_SPECS="${1:-tests/e2e/*.spec.ts}"

echo "==> Verification de la disponibilite de la stack Docker Odoo 18..."
docker-compose -f "$COMPOSE_FILE" up -d

echo "==> Attente active (polling HTTP) du serveur Odoo sur le port 8069..."
until curl -s -f -o /dev/null "http://localhost:8069/web/login"; do
  sleep 1
done

echo "==> Lancement de la campagne de tests de navigation Playwright E2E..."
npx playwright test $TARGET_SPECS --reporter=list,html

echo "==> Rapport HTML et videos generes dans artifacts/playwright-report/"
\end{lstlisting}